\clearpage
\section{Input: three paths to the same methods}

There are exactly three ways to reach the timer: keyboard, menu bar, and buttons. All
three end at the same \texttt{App} methods.

\subsection{Keyboard}

\begin{figure}[H]
\centering
\begin{tikzpicture}[
  k/.style={rounded corners=2pt,draw=ink!50,fill=white,align=center,inner sep=4pt,
            font=\scriptsize\sffamily,text width=30mm,minimum height=9mm},
  g/.style={k,draw=mint!70!black,fill=mint!18,text=mint!60!black},
  a/.style={k,draw=accent!65,fill=accent!12,text=accent!55!black},
]
  \node[k] (sp) {\texttt{Space}};
  \node[k] (r) {\texttt{R}};
  \node[k] (s) {\texttt{S}};
  \node[k] (x) {\texttt{X}};
  \node[a] (f) {\texttt{F}};
  \node[k] (esc) {\texttt{Escape}};

  \node[g] (t1) at (5.6,1.6) {\texttt{toggle()}};
  \node[g] (t2) at (5.6,0.4) {\texttt{restart()}};
  \node[g] (t3) at (5.6,-0.8) {\texttt{skip()}};
  \node[g] (t4) at (5.6,-2.0) {\texttt{stop()}};
  \node[g] (t5) at (5.6,-3.2) {\texttt{complete\_current()}\\\scriptsize only if \texttt{is\_overtime()}};
  \node[g] (t6) at (5.6,-4.4) {\texttt{panel = None}\\\scriptsize closes any open panel};

  \draw[e] (sp)--(t1); \draw[e] (r)--(t2); \draw[e] (s)--(t3);
  \draw[e] (x)--(t4); \draw[e] (f)--(t5); \draw[e] (esc)--(t6);
\end{tikzpicture}
\caption{Six keys. \texttt{Escape} is handled first and closes the panel; the other five
are ignored while a panel is open or while egui wants the keyboard.}
\label{fig:keys}
\end{figure}

\begin{lstlisting}[caption={\file{app.rs} lines 190--216. The guard at the top is what
keeps the timer from reacting while the user is typing in a settings field.}]
pub fn handle_keys(&mut self, ctx: &egui::Context) {
    let escape = ctx.input(|i| i.key_pressed(egui::Key::Escape));
    if escape { self.panel = Panel::None; self.confirm_reset = false; }
    if ctx.egui_wants_keyboard_input() || self.panel != Panel::None { return; }
    ctx.input(|i| {
        if i.key_pressed(egui::Key::Space) { self.toggle(); }
        if i.key_pressed(egui::Key::R)     { self.restart(); }
        if i.key_pressed(egui::Key::S)     { self.skip(); }
        if i.key_pressed(egui::Key::X)     { self.stop(); }
        if i.key_pressed(egui::Key::F) && self.is_overtime() { self.complete_current(); }
    });
}
\end{lstlisting}

\subsection{Global hotkeys, without a keyboard tap}

The AppKit shell registers five Carbon hotkeys. They work even when the app is not
focused, and they need no permission prompt because they are registered on the
application's own event target.

\begin{figure}[H]
\centering
\begin{tikzpicture}[
  b/.style={rounded corners=2pt,draw=ink!50,fill=white,align=center,inner sep=4pt,
            font=\scriptsize\sffamily,text width=40mm,minimum height=10mm},
  k/.style={b,draw=accent!65,fill=accent!12,text=accent!55!black},
  g/.style={b,draw=mint!70!black,fill=mint!18,text=mint!60!black},
]
  \node[b] (a) {\textbf{RegisterEventHotKey}\\\scriptsize Control + Option};
  \node[b] (b) {\textbf{key codes}\\\texttt{49} Space\quad \texttt{15} R\quad \texttt{1} S\\\texttt{7} X\quad \texttt{17} T};
  \node[b] (c) {\textbf{InstallEventHandler}\\\scriptsize on \texttt{GetApplicationEventTarget()}};
  \node[k] (d) {\textbf{callback}\\\texttt{key\_id(event)}\\\scriptsize checks the \texttt{TMRs} signature};
  \node[g] (e) {\textbf{dispatch}\\\texttt{0} toggle\quad \texttt{1} restart\\\texttt{2} skip\quad \texttt{3} stop\\\texttt{4} show};
  \draw[e] (a)--(b); \draw[e] (b)--(c); \draw[e] (c)--(d); \draw[e] (d)--(e);
\end{tikzpicture}
\caption{\file{hotkeys.rs} is 128 lines of raw Carbon FFI. The \texttt{TMRs} signature
is a magic value that stops the handler from reacting to some other app's hotkey.}
\label{fig:hotkeys}
\end{figure}

\begin{table}[H]
\centering\small
\begin{tabular}{@{}p{22mm} p{26mm} p{30mm} p{60mm}@{}}
\toprule
\textbf{Key code} & \textbf{Key} & \textbf{Modifiers} & \textbf{Action} \\
\midrule
49 & Space & Control + Option & \texttt{toggle()} \\
15 & R & Control + Option & \texttt{restart()} \\
1 & S & Control + Option & \texttt{skip()} \\
7 & X & Control + Option & \texttt{stop()} \\
17 & T & Control + Option & \texttt{show()} \\
\bottomrule
\end{tabular}
\caption{The five registered hotkeys. The \texttt{T} key is the ``show timer'' command;
it is the only one that does not mutate the timer.}
\end{table}

\subsection{The menu bar}

\begin{figure}[H]
\centering
\begin{tikzpicture}[
  m/.style={rounded corners=2pt,draw=ink!50,fill=white,align=center,inner sep=3pt,
            font=\scriptsize\sffamily,text width=34mm,minimum height=8mm},
  h/.style={m,draw=accent!65,fill=accent!12,text=accent!55!black},
]
  \node[m] (i0) {\texttt{Start / Pause}};
  \node[m] (i1) {\texttt{Restart}};
  \node[m] (i2) {\texttt{Skip}};
  \node[m] (i3) {\texttt{Stop}};
  \node[m] (i4) {\texttt{Finish overtime}};
  \node[m] (i5) {\texttt{Show timer}};
  \node[m] (i6) {\texttt{Settings...}};
  \node[m] (i7) {\texttt{Statistics...}};
  \node[m] (i8) {\texttt{Reset session counter}};
  \node[h] (i9) {\texttt{Quit}};
  \foreach \n in {i0,i1,i2,i3,i4,i5,i6,i7,i8,i9}
    \draw[ethin, softink!35] (\n.south) -- ++(0,-3mm);
  \node[note, below=5mm of i4, text width=110mm] {%
    The same ten items exist in both shells. In the egui shell they are produced by
    \file{native.rs}; in the AppKit shell by \file{mac\_shell.rs}. The tags are identical,
    so the \texttt{match} in \texttt{act:} and the \texttt{match} in \texttt{logic()}
    agree by construction.};
\end{tikzpicture}
\caption{Ten menu items. The first item's title tracks the run state, so the menu bar
always says \texttt{Start}, \texttt{Pause}, or \texttt{Resume}.}
\label{fig:menubar}
\end{figure}